\documentclass[11pt,a4paper]{article}
\usepackage[T1]{fontenc}
\usepackage[utf8]{inputenc}
\usepackage{amsmath,amssymb,xcolor}
\usepackage[margin=25mm]{geometry}
\usepackage[colorlinks=true,urlcolor=blue]{hyperref}
\definecolor{nexorablue}{HTML}{2459B5}
\setlength{\parindent}{0pt}
\setlength{\parskip}{7pt}
\setlength{\emergencystretch}{2em}
\renewcommand{\familydefault}{\sfdefault}
\begin{document}

{\small\bfseries\color{nexorablue} NEXORA CPGE}\par
{\small MPSI / NEX-RGT-03 / Énoncé}\par
{\LARGE\bfseries Charge depuis une tension non nulle}\par
\vspace{0.25em}\hrule\vspace{0.6em}

Une capacité de $47\,\mathrm{\mu F}$ est initialement portée à $2{,}0\,\mathrm V$. On la relie à une source de $6{,}0\,\mathrm V$ via $10\,\mathrm{k\Omega}$. Établir $u_C(t)$, calculer le courant initial et l’instant où $u_C$ atteint $5{,}0\,\mathrm V$.


\end{document}
